\section{Tests and additional results}
\label{app:extra}

\subsection{Losses and tests of the pre-registered comparisons}
\label{app:tests}

\begin{table}[!htb]
\centering
\footnotesize
\setlength{\tabcolsep}{2.0pt}
% Generated by scripts/paper_results.py from runs/. Do not edit.
\begin{tabular}{@{}lllcccccc@{}}
\toprule
& & Training & \multicolumn{2}{c}{Loss} & & & & \\
\cmidrule(lr){4-5}
Ours & Reference & tokens & ours & ref. & $\Delta$ & $\Delta<0$ & $p_{\mathrm{Holm}}$ & $\Delta$, cache \\
\midrule
\multicolumn{9}{@{}l}{\emph{GDN with \ingram{} and no softmax attention, against 3:1 hybrids}} \\
$8\times256$ (56.0M) & $8\times256$ (30.4M) & 164M & $0.607$ & $0.749$ & $-0.142$ & 8/8 & $0.012$ & $-0.046$ \\
 &  &  &  &  & {\scriptsize [$-0.166$, $-0.127$]} &  &  & {\scriptsize [$-0.054$, $-0.042$]} \\
$8\times256$ (56.0M) & $12\times384$ (54.2M) & 164M & $0.607$ & $0.709$ & $-0.102$ & 8/8 & $0.008$ & $-0.023$ \\
 &  &  &  &  & {\scriptsize [$-0.113$, $-0.092$]} &  &  & {\scriptsize [$-0.028$, $-0.018$]} \\
$8\times256$ (56.0M) & $8\times256$ + table (55.8M) & 164M & $0.607$ & $0.745$ & $-0.137$ & 8/8 & $0.016$ & $-0.042$ \\
 &  &  &  &  & {\scriptsize [$-0.179$, $-0.104$]} &  &  & {\scriptsize [$-0.056$, $-0.030$]} \\
$12\times384$ (80.3M) & $12\times384$ (54.2M) & 164M & $0.582$ & $0.709$ & $-0.127$ & 8/8 & $0.016$ & $-0.044$ \\
 &  &  &  &  & {\scriptsize [$-0.137$, $-0.118$]} &  &  & {\scriptsize [$-0.049$, $-0.040$]} \\
$16\times512$ (115.3M) & $16\times512$ (88.4M) & 328M & $0.503$ & $0.561$ & $-0.058$ & 4/4 & n/a & $-0.021$ \\
 &  &  &  &  & {\scriptsize [$-0.080$, $-0.040$]} &  &  & {\scriptsize [$-0.033$, $-0.012$]} \\
$8\times256$ (56.0M) & $8\times256$ (30.4M) & 328M & $0.550$ & $0.622$ & $-0.072$ & 8/8 & $0.016$ & $-0.024$ \\
 &  &  &  &  & {\scriptsize [$-0.085$, $-0.061$]} &  &  & {\scriptsize [$-0.031$, $-0.018$]} \\
$8\times256$ (56.0M) & $8\times256$ (30.4M) & 164M/328M & $0.607$ & $0.622$ & $-0.015$ & 6/8 & $0.039$ & $+0.026$ \\
 &  &  &  &  & {\scriptsize [$-0.028$, $-0.003$]} &  &  & {\scriptsize [$+0.020$, $+0.032$]} \\
\addlinespace
\multicolumn{9}{@{}l}{\emph{7:1 hybrid with \ingram{}, against hybrids without it}} \\
$8\times256$ (55.9M) & 3:1, $8\times256$ (30.4M) & 1.31B & $0.448$ & $0.510$ & $-0.063$ & 8/8 & $0.008$ & $-0.036$ \\
 &  &  &  &  & {\scriptsize [$-0.071$, $-0.054$]} &  &  & {\scriptsize [$-0.042$, $-0.031$]} \\
$8\times256$ (55.9M) & 7:1, $8\times256$ (30.4M) & 1.31B & $0.448$ & $0.508$ & $-0.060$ & 8/8 & $0.008$ & $-0.034$ \\
 &  &  &  &  & {\scriptsize [$-0.067$, $-0.055$]} &  &  & {\scriptsize [$-0.039$, $-0.031$]} \\
\bottomrule
\end{tabular}

\caption{Losses and tests for the comparisons of Figure~\ref{fig:confirm}, in the same order, at
16k context on seeds 20--27 ($n = 8$, paired; $n = 4$ at $16\times512$). Models are given as
blocks $\times$ width with their parameter count; in the first group every reference is a 3:1
hybrid, and ``+ table'' marks the one with a static $N$-gram table. Loss is the mean over seeds; $\Delta$ is ours
minus reference, with its 95\% bootstrap interval over seeds below it. $p_{\mathrm{Holm}}$ is the
one-sided Wilcoxon $p$ after Holm correction within the pre-registered family. In the row marked
164M/328M our model trains for 164M tokens and the reference for 328M
(Section~\ref{sec:length}). The with-cache difference was tested in the first two rows (same
adjusted $p$) and is reported without a test in the others. Per-seed losses are in
Tables~\ref{tab:seeds164} and~\ref{tab:seedslong}.}
\label{tab:confirm}
\end{table}

Table~\ref{tab:confirm} lists the losses and tests behind Figure~\ref{fig:confirm}. Across the 32
evaluation windows, the share of tokens that the lookup of Section~\ref{sec:repeats} predicts
ranges from 63.4\% to 98.4\%, and from 8k tokens of depth on, tokens for which an older
continuation is the right one stay between 5.3\% and 7.9\% (Figure~\ref{fig:repeats}b). In the
first row, the mean difference without the outlier seed is $-0.131$. The hybrid of the
parameter-matched control has 54.2M parameters, of which 15.6M lie outside the embedding and
output layers,
against 5.1M for our model outside embedding, output and table. With 12 blocks of width 384, GDN
with \ingram{} has 80.3M parameters against 54.2M for the 3:1 hybrid, and with 16 blocks of width
512, trained for 20,000 steps on four seeds per arm, 115.3M against 88.4M.

\subsection{Training length}
\label{app:length}

Between 164M and 328M training tokens the hybrid improves by 0.127 nats and the model with
\ingram{} by 0.057. On seeds 0 and 1 at 164M tokens, the 7:1 hybrid reaches 0.995 without
\ingram{} and 0.607 with it, against 0.722 for the 3:1 hybrid (Figure~\ref{fig:grid}); the
1.31B-token runs test whether this ordering persists with eight times more training. Without
\ingram{}, the benefit of the second attention layer decreases with training length (post hoc). It
is 0.27 nats at 164M tokens on seeds 0 and 1, and 0.075 nats at 8,000 steps of the long runs. From
56,000 steps on, the two means differ by at most 0.003 (Figure~\ref{fig:training}a), and at the
end 3 of 8 seeds have the 3:1 hybrid below the 7:1 hybrid. The model with \ingram{} also varies
less across seeds at the end of training: the standard deviation over seeds is 0.002, against
0.013 for the 3:1 hybrid and 0.010 for the 7:1 hybrid.

\subsection{Step time and state}
\label{sec:cost}

\begin{table}[!htb]
\centering
\footnotesize
% Generated by scripts/paper_results.py from runs/. Do not edit.
\begin{tabular}{@{}lcccccc@{}}
\toprule
& & & \multicolumn{2}{c}{Seconds per step} & Peak GPU & State at 16k \\
\cmidrule(lr){4-5}
Model & Steps & Parameters & median & range & memory (GB) & tokens (MB) \\
\midrule
3:1 hybrid, $8\times256$ & 10k & 30.4M & 0.122 & 0.101--0.125 & 4.3 & 67.5 \\
GDN + \ingram{}, $8\times256$ & 10k & 56.0M & 0.130 & 0.116--0.198 & 5.8 & 12.7 \\
3:1 hybrid + static table, $8\times256$ & 10k & 55.8M & 0.138 & 0.123--0.148 & 5.7 & 67.5 \\
3:1 hybrid, $12\times384$ & 10k & 54.2M & 0.156 & 0.150--0.175 & 6.6 & 151.9 \\
GDN + \ingram{}, $12\times384$ & 10k & 80.3M & 0.167 & 0.161--0.185 & 8.2 & 13.3 \\
\midrule
3:1 hybrid, $8\times256$ & 80k & 30.4M & 0.122 & 0.099--0.135 & 4.3 & 67.5 \\
7:1 hybrid, $8\times256$ & 80k & 30.4M & 0.122 & 0.093--0.127 & 4.4 & 34.0 \\
7:1 hybrid + \ingram{}, $8\times256$ & 80k & 55.9M & 0.131 & 0.118--0.160 & 5.8 & 46.2 \\
\bottomrule
\end{tabular}

\caption{Cost of the models in Table~\ref{tab:confirm}. Seconds per step are wall-clock time per
training step, including data loading and the periodic evaluations, over seeds 20--27; at each
evaluation a model with \ingram{} runs three additional passes with its reads ablated, which are
included. Each run used one rented RTX~4090, and hosts differ between runs. State is the fp32 size of what a decoder
stores per session after 16,384 tokens, computed as described in Appendix~\ref{sec:cost};
Figure~\ref{fig:state} extends it to longer contexts.}
\label{tab:cost}
\end{table}

Training steps with \ingram{} were slower in our runs (Table~\ref{tab:cost}). Over seeds 20--27,
GDN with \ingram{} takes a median of 0.130 seconds per step (range 0.116 to 0.198) and the 3:1
hybrid 0.122 (0.101 to 0.125). Because the runs used different rented hosts, the times are
unpaired: the ratio of medians is 1.07 and the ratio of means is 1.19. The test of
Section~\ref{sec:length} that allows the hybrid twice the steps favours the hybrid for any
step-time ratio below 2.

We compute the state that a decoder stores per session at 16,384 tokens, in fp32
(Table~\ref{tab:cost}): keys and values of the model width for every token in each attention
layer, the recurrent state of each GDN block, and one 64-dimensional vector per written slot of
\ingram{} (47,426 slots on average for a held-out session read from its first token; the hash
index is not counted). This amounts to 12.7~MB for GDN with \ingram{}, 67.5~MB for the 3:1 hybrid
and 46.2~MB for the 7:1 hybrid with \ingram{}. The evaluation windows are cut inside sessions and
contain fewer new cues: 29,337 slots on average, or 1.79 per token.

\subsection{Dose grid and attention statistics}
\label{app:grid}

In the grid of Figure~\ref{fig:grid}, four attention layers reach 0.732 at their better learning
rate, and a model of eight attention layers with no GDN reaches 0.766 at its better learning rate
(Table~\ref{tab:lr}). On these two seeds the 3:1 hybrid is therefore the best model without
\ingram{}. With two placements the four cells lie between 0.584 and 0.590.

The summary we fixed in advance is the iso-loss attention $A^{*}(P)$: the number of attention
layers that a model with $P$ placements needs in order to match the 3:1 hybrid without \ingram{}
(0.722), found by linear interpolation over $A$. For $P = 1$, $2$ and $4$, the cell with no
attention is already below that loss on both seeds, by 0.113, 0.131 and 0.138 nats in the seed
mean, and hence $A^{*}(P) \le 0$ for every $P$ we ran.

With four placements the cost of removing both attention layers is 0.001 over all tokens
($[-0.001, 0.003]$), and the intervals for older-value tokens ($+0.016$, $[0.003, 0.029]$) and for
tokens with no match ($+0.013$, $[0.005, 0.021]$) still exclude zero (Table~\ref{tab:a3}).
\ingram{} also covers near-duplicate lines, which we had expected to depend on attention: their
cost decreases from $+0.261$ at $P = 0$ to $+0.011$ at $P = 1$. Tokens with no match change least
($+0.048$ and $+0.053$). On seeds 20--27 without attention, the model with \ingram{} is also worse
than the 3:1 hybrid on format tokens ($+0.012$, $[0.005, 0.020]$ at 328M tokens); with one
attention layer, at 1.31B tokens, it is not ($-0.007$, $[-0.013, -0.004]$).

To measure where the attention layers attend, we define the induction share of far attention as
the fraction of the attention mass on keys more than 512 tokens back that falls on positions
directly after an earlier occurrence of the current bigram, averaged over heads and layers
(Table~\ref{tab:a4}). In the 3:1 cells it is 0.233 without \ingram{} and 0.136, 0.181 and 0.174 at
$P = 1$, $2$ and $4$; with four attention layers it is 0.211 and then 0.154, 0.150 and 0.109. The
share is lower with \ingram{} in the seed mean of all six cells and in ten of the twelve
single-seed comparisons.

\subsection{Cue lengths and further checks}
\label{sec:cues}

On seeds 0 and 1, longer cues lower the loss (Table~\ref{tab:cues}): 0.633 with cue lengths 1 to
5, 0.614 with $\{1, 2, 4, 8, 16\}$ at the same parameter count, and 0.606 with $\{1, 2, 4, 8, 16,
32\}$. We selected between $\{1, 2, 3, 5, 8, 16\}$ and $\{1, 2, 4, 8, 16, 32\}$ on these seeds by
a rule fixed in advance, the lower mean loss. On seeds 20--27, the selected set reaches 0.607 against 0.632 for cue lengths 1
to 5, $\Delta = -0.025$ ($[-0.028, -0.022]$; 8 of 8; Holm $p = 0.012$); this was the primary test
of its family. The n-gram cache contributes less when the cues are long: it lowers the loss by
0.038 nats with cue lengths 1 to 5 and by 0.027 with the selected set, and with the cache on both
sides the difference between the two cue sets is $-0.014$ ($[-0.016, -0.012]$; reported, not
tested). Placements and long cues combine: on seeds 0 and 1, two placements reach 0.608 with cue
lengths 1 to 5 and 0.590 with the selected set (Table~\ref{tab:grid}).

Appendix~\ref{app:other} reports three further checks: a scored read rule lowers the loss by
0.005 nats in a smaller setting and was worse than the additive read on both selection seeds of
this recipe; rounding the rows of the hashed table and the slot state at each read to 4 bits
changes the loss by at most 0.0005 nats; and we detect no increase in looping under generation at 32 prompts and two seeds.

\subsection{Slot growth and the serving scenarios}
\label{app:serving}

The slot rate of Section~\ref{sec:serving-kind} varies between sessions: after 16,384 tokens it
ranges from 1.19 to 4.06 slots per token across the 704 held-out trajectories, and after 131,072
tokens from 0.55 to 2.02 across the 134 trajectories of that length. One held-out trajectory
reaches 262,144 tokens; at positions between 128k and 256k, a new token adds 0.55 slots on average
over the trajectories that reach them. Long cues account for most of the slots: after 16,384
tokens, cue length 32 has written 0.82 slots per token and cue length 1 has written 0.10. Our decoder stores the
hash index on the host and the slot vectors on the accelerator, and all three of our model sizes
use slot width 64. Our implementation also has a bounded store, in which a cue hashes to a set of
four slots and a new cue evicts the least recently used slot of its set; a device with a fixed
memory budget requires such a store, and we plan to evaluate it on this recipe.

In the scenario of Table~\ref{tab:serving}a, each attention layer requires 2.1~KB per token and
\ingram{} 1.6~KB per token in total; the added context comes from the attention caches that are
removed. At the bound of six slots per token, the largest contexts on the laptop are 559k tokens
with six attention layers and 719k with four. Decode speed is unchanged in the roofline model (142
to 146 tokens per second at 262,144 tokens of context on the laptop), because sparse attention
reads a bounded number of tokens per step. In a decode benchmark of an earlier configuration,
\ingram{} added 1.3 to 1.4~ms per token to GDN, and neither model slowed down between 4,096 and 1M
tokens of context (Appendix~\ref{app:decode}).

